% ID: SEL-D0079-OB-E01
% Tanda S079.
\textbf{a)} En $(7,3)$ se cumple $3\cdot7+4\cdot3=33\ge28$, $5\cdot7+2\cdot3=41\le42$ y $7-3=4\ge0$. Por tanto, el punto pertenece al recinto.

\textbf{b)} Despejamos $y$:
\[
 y\ge7-\frac34x,\qquad y\le21-\frac52x,\qquad y\le x.
\]
El recinto es el triángulo cerrado situado por encima de la primera recta y por debajo de las otras dos. Intersecando cada pareja de fronteras, obtenemos
\[
 \begin{array}{c|c}
 \text{Fronteras}&\text{Vértice}\\\hline
 3x+4y=28,\ y=x&(4,4)\\
 5x+2y=42,\ y=x&(6,6)\\
 3x+4y=28,\ 5x+2y=42&(8,1)
 \end{array}
\]
Todos satisfacen la tercera restricción de cada pareja.
\begin{center}
\setlength{\unitlength}{1mm}
\begin{picture}(102,70)(-9,-9)
\put(0,0){\vector(1,0){86}}\put(0,0){\vector(0,1){58}}
\put(88,0){$x$}\put(1,58){$y$}
{\color{gray}
\qbezier(0,49)(33.6,24.5)(67.2,0)
\qbezier(38.88,52.5)(49.68,26.25)(60.48,0)
\qbezier(0,0)(28.8,28)(57.6,56)}
\linethickness{0.45mm}
\qbezier(28.8,28)(36,35)(43.2,42)
\qbezier(43.2,42)(50.4,24.5)(57.6,7)
\qbezier(57.6,7)(43.2,17.5)(28.8,28)
\put(17,28){\scriptsize $(4,4)$}
\put(43,45){\scriptsize $(6,6)$}
\put(59,7){\scriptsize $(8,1)$}
\end{picture}
\end{center}

\textbf{c)} Evaluamos la función lineal en los vértices del polígono cerrado y acotado:
\[
 F(4,4)=10,\qquad F(6,6)=12,\qquad F(8,1)=28.
\]
El máximo es $\boxed{28}$ y se alcanza únicamente en $(8,1)$.
